\documentclass[11pt]{article}
\usepackage[margin=2.2cm]{geometry}
\usepackage{amsmath,amssymb,listings}
\usepackage{booktabs}
\usepackage{xurl}
\usepackage{xcolor}
\usepackage[colorlinks=true,linkcolor=blue!50!black,citecolor=blue!50!black,urlcolor=blue!50!black]{hyperref}
\usepackage{orcidlink}

\lstset{
    basicstyle=\ttfamily\footnotesize,
    keywordstyle=\color{blue!60!black}\bfseries,
    stringstyle=\color{green!45!black},
    commentstyle=\color{gray}\itshape,
    breaklines=true,
    frame=single,
    framerule=0.3pt,
    xleftmargin=1em,
}
\lstdefinelanguage{Rust}{
  morekeywords={fn,let,mut,if,else,match,return,while,for,in,pub,struct,impl,u32,u64,usize,bool,use},
  sensitive=true,
  morecomment=[l]{//},
  morestring=[b]",
}
\input{gen/rustproof_facts.tex}

\title{Proof-Carrying Rust\\
\large Logic safety on top of memory safety, in a toolchain small enough to read}
\author{Brent S. Hartshorn \orcidlink{0009-0004-2853-655X}
        \small (\url{brenthartshorn@proton.me})
\\
\small
\url{https://github.com/brentharts/RosettaMath}
\quad
\url{https://github.com/brentharts/crust}
}
\date{September 22, 2026}

\begin{document}
\maketitle

\begin{abstract}
Rust's guarantee ends at memory.  A function that returns the wrong number
is well-typed Rust, and in the domains that still run on OCaml a wrong
number has a price.  Crust, a Rust front end written in Python inside a C
toolchain, now reads contracts in the syntax Creusot and Prusti already use
--- \texttt{\#[requires]}, \texttt{\#[ensures]}, \texttt{\#[invariant]},
\texttt{\#[variant]} --- so an annotated file stays ordinary Rust.  The
compiler checks every clause at runtime by default.  A source-level lift
turns the same function and the same clauses into the fragment a proof
kernel of \kernelLines{} lines reads, and a second lift states every place
the function could panic --- an overflow, an underflow, a division, an index,
a broken precondition --- as a theorem of its own.  Every proof the kernel
accepts is exported to Lean~4 and checked again.  On the LeanOS register
sizer written as a Rust \texttt{match}, both postconditions are proved for
every input, a false bound is refused, and Lean \regsLeanVerdict{} the export
with \regsLeanNoAxioms{} theorem depending on no axioms.  On the LeanOS
allocator ported to Rust, \allocProved{} of \allocObligations{} panic
obligations are proved from the shipped source --- after three the prover
left open turned out to be real overflows, which the runtime contract check
confirmed and the source now avoids.  Every one of the \rustCertificates{} Rust theorems
the kernel settles is exported, and Lean accepts \rustLeanAgreed{} of them
with no axioms.  Along the way the kernel's unary numerals, which
made \texttt{u32::MAX} four billion nested terms, became one node each.  We
argue for two checkers rather than a better one: after the 2026 soundness
bugs in Lean's kernel, a proof is worth more when a second, independently
written kernel agrees.
\end{abstract}

\section{Introduction: the gap Rust leaves}

Rust made a class of bug a compile-time matter.  Use after free, data
races, buffer overruns in safe code: the borrow checker refuses them before
the program runs, and that is why Rust is arriving everywhere systems
software is written.  But the guarantee is about memory.  A function that
takes a class number and returns the wrong register count is well-typed
Rust; so is one that computes a price, a margin, a settlement date, and
gets it wrong.  Rust has no more to say about logic than C did.

The domains where logic errors cost the most have noticed.  Trading firms
kept OCaml for decades, in part because its type system and its culture of
small, pure, reasoned-about functions let a team believe a function does
what it says~\cite{biggar2020}.  Moving to Rust means gaining memory safety
and losing nothing at the type level, but it does not gain a way to state
what a function is \emph{for} and have the toolchain hold it to that.

This paper closes that gap for a dialect of Rust, and the closing has four
parts, each of which is already in the shipped toolchain:

\begin{enumerate}
\item Contracts written as Creusot and Prusti write them, so the file is
  still Rust for those tools and for \texttt{rustc}
  (Section~\ref{sec:contracts}).
\item Every clause checked at runtime, at every exit, by default
  (Section~\ref{sec:contracts}).
\item A lift from the Rust source, not from an intermediate language, into
  the fragment a small proof kernel reads; and a second lift that turns every
  potential panic into a theorem (Sections~\ref{sec:lift}
  and~\ref{sec:safety}).
\item Two checkers.  The kernel that proves is \kernelLines{} lines of
  Python; Lean~4 re-checks what it accepts (Section~\ref{sec:two}).
\end{enumerate}

The claim is not that this is the only Rust that can be verified.  Aeneas,
Hax, Creusot and Prusti exist, and Klaus, Conejero and Tolmach report a
working pipeline from production Rust to Lean~4 proofs closed by AI
provers~\cite{klaus2026}.  The claim is narrower and, we think, more useful:
that contracts, runtime checks, source-level proof and independent
re-checking are one compile, in one toolchain a person can read end to
end, with every number in this paper produced by the run that typeset it.

\section{Contracts that are still Rust}\label{sec:contracts}

Crust accepts contract attributes in the syntax of Creusot and Prusti
\cite{creusot,prusti}, including their crate paths, and \texttt{use
prusti\_contracts::*;} compiles to nothing.  A file annotated for Crust is
valid input to those tools and, with their shim crates, to \texttt{rustc}.

\begin{lstlisting}[language=Rust]
#[requires(x < 100)]
#[ensures(result <= 23)]
fn regs(c: u32) -> u32 { match c { 0 => 6, 1 => 10, 2 => 15, _ => 23 } }

#[ensures(*v == old(*v) + 1)]
fn inc(v: &mut u32) { *v += 1; }

#[invariant(i <= n)]
#[variant(n - i)]
while i < n { s += i; i += 1; }
\end{lstlisting}

The compiler lowers each clause to a runtime check: \texttt{requires} at
entry; \texttt{ensures} at every exit, which includes an explicit
\texttt{return}, a tail expression, each arm of a tail \texttt{match}, the
early return of \texttt{?}, and falling off the end of a function that
returns nothing; \texttt{invariant} at every loop head, including the head
whose test fails; \texttt{variant} as a strict decrease.  Inside a clause,
\texttt{result} is the return value, \texttt{old(e)} is \texttt{e}
snapshotted at entry, and \texttt{a ==> b} is implication.  A violation
prints the clause as written, the function and the line, then aborts:

\begin{lstlisting}
contract violated: ensures `result <= 23` of `regs` (line 2)
\end{lstlisting}

Clauses with \texttt{forall} or \texttt{exists} are kept for the prover and
not run.  \texttt{CRUST\_CONTRACTS=0} turns the checks off, and the clauses
are still parsed, so a malformed one is an error either way.  The
\contractTests{} tests of this feature each compile a program, run it, and
read the exit status and the message.

Runtime checking is not verification, and it comes first on purpose.  It is
what the code does today, in every build, with no proof tooling installed;
and it keeps a discipline the earlier papers in this series relied on: a
contract the test corpus breaks is a contract no proof will
establish~\cite{hartshorn2026lean}.

\section{From the source, not the IL}\label{sec:lift}

The LeanOS papers reached compiled code through the intermediate language:
\texttt{ilproof.py} lifts a function from the compiler's IL back into the
imperative fragment that \texttt{hoare.py} proves things
about~\cite{hartshorn2026leanos}.  That works for C.  For idiomatic Rust it
reaches almost nothing, and the reason is instructive: the IL has already
thrown away what a proof needs.  A \texttt{match} arrives as a
\texttt{switch}, which the lift sees as \texttt{goto}s.  A tail
\texttt{if}/\texttt{else} leaves a dangling implicit return.  \texttt{\&\&}
is a pair of jumps.  And a \texttt{u32} is an \texttt{int} of no particular
width, which is exactly the information an overflow obligation is made of.

\texttt{shivyc/rustproof.py} (\liftLines{} lines) lifts from the source
instead, where all of that is still present, and it reads the contracts on
the function.  \texttt{leanos/regs.rs} is the register sizer of
LeanOS~\cite{hartshorn2026leanos} written as Rust writes it:

\lstinputlisting[language=Rust,firstline=20,lastline=29]{\crustDir/leanos/regs.rs}

It lifts to what a person would write, the same fragment the IL lift
produces from the C version in \texttt{crustos/kernel.c}, with the
postconditions beside it:

\lstinputlisting[language=Python]{gen/elf_regs_lifted.py}

\texttt{by\_every\_bool} --- split on every guard, then compute --- proves
\texttt{\regsEnsures} for every \texttt{cls} in \regsSeconds{}\,s, and the
false bound \texttt{result <= 15} is \regsFalseBound.  The Rust model and the
C model agree on every class.  This is the first function in the tree whose
model and theorem are both read off the file that ships; the LeanOS paper
spent a section on the transcription problem --- a model re-typed by hand
beside the code, and tests to keep the two aligned --- and here there is
nothing to transcribe.

What the lift reaches: unsigned integers (as Nat) and \texttt{bool};
slices and \texttt{\&Vec} as arrays; structs declared in the file as records,
with one \texttt{\&mut} struct parameter threaded as state; \texttt{let} and
assignment; \texttt{if} and \texttt{match} as statements, values or tails,
\texttt{match} with literals, alternatives, ranges, bindings and guards;
early \texttt{return}; \texttt{for} over ranges; \texttt{while} with
\texttt{\#[invariant]} and \texttt{\#[variant]}; and calls within the file.
Everything else --- signed integers, methods, recursion, bitwise operators,
narrowing casts, quantified clauses --- is refused with the construct named.
The refusals are the roadmap, as they were for the IL lift.

\section{Every panic is a theorem}\label{sec:safety}

A theorem about the lifted model is a theorem about Nat, and Rust's
integers are not Nat: \texttt{a + b} wraps or panics past the type's
maximum, and \texttt{a - b} panics below zero where Nat truncates.  Rather
than weaken the model, the lift states the difference as obligations.
\texttt{Lifted.safety()} is a second lift of the same function,
\texttt{f\_\_safe}, with the same control flow, in which a boolean
\texttt{\_ok} is conjoined with a claim at each place the Rust would panic
and every \texttt{return} returns \texttt{\_ok}.  The theorem
``\texttt{f\_\_safe(args)} for every \texttt{args}'' is ``no input reaches
a panic''.

\begin{center}
\begin{tabular}{ll}
\toprule
Rust & owes \\
\midrule
\texttt{a + b}, \texttt{a * b}, \texttt{a << k} in a \texttt{uN} & the value $\le$ \texttt{max\_uN} \\
\texttt{a - b} & \texttt{b <= a} \\
\texttt{a / b}, \texttt{a \% b} & \texttt{b != 0} \\
\texttt{xs[i]} & \texttt{i < xs.len()} \\
\texttt{f(args)} & \texttt{f}'s \texttt{\#[requires]} of \texttt{args} \\
\bottomrule
\end{tabular}
\end{center}

Each obligation is paid before the statement that evaluates it, so it reads
the variables as they are then, and one the Rust evaluates only on some
paths --- the right of \texttt{\&\&}, a \texttt{match} guard --- is
conditioned on the path.  A \texttt{\#[requires]} guards the body.  For a
function that returns early when a subtraction would underflow:

\lstinputlisting[language=Python]{gen/early_sub_safe.py}

\texttt{tools/rustprove.py} (\proveLines{} lines) proves each obligation on
its own and reports it \emph{proved} or \emph{open}.  Open means the
automation did not settle it and says nothing about whether it holds.  Over
\safetyShapes{} small functions chosen to be safe or unsafe by construction,
the verdicts \safetyAgrees{} with the construction:

\begin{center}
\begin{tabular}{lrl}
\toprule
function & obligations & verdict \\
\midrule
\input{gen/safety_table.tex}
\end{tabular}
\end{center}

Three adjustments to the kernel's tactics were needed to get there, each
checked by the kernel rather than trusted.  An obligation's hypotheses are
weakened away and the proof re-wrapped, since \texttt{by\_every\_bool}
does not bind them.  An obligation is stated in every spelling a guard
might use, because \texttt{if i >= len \{ return \}} decides
\texttt{len <= i} and not \texttt{i < len}, which over Nat is the same fact
and as a boolean is a different term.  And \texttt{if a \&\& b} with no
\texttt{else} is lowered as nested \texttt{if}s, so that each conjunct is a
guard the split can decide.

\subsection{The allocator, ported and proved}

\texttt{leanos/alloc.rs} is the LeanOS per-thread bump
allocator~\cite{hartshorn2026leanos}, ported from RPython to Rust so that
no input can make it panic --- every array read behind its own length
check, \texttt{tid + 1} rewritten so it cannot overflow, every addition
checked before it is made (Section~\ref{sec:bounds}) --- and checked
against the Python row for row on the \allocCorpusRows-row corpus of the
original model test.  All \allocFunctions{} functions lift
(\allocRefused{} refused).  \texttt{bump}, the function that moves the
counter:

\lstinputlisting[language=Python]{gen/bump_lifted.py}

Its \texttt{\#[ensures(used <= result)]} is the theorem the LeanOS paper
called \texttt{bump\_monotone}, there proved about a hand-typed model; here
it is \bumpMonotone{} about the generated one, and \texttt{bump\_bounded}
is proved against the same generated model.  Of
the \allocObligations{} panic obligations in the file, \allocProved{} are
proved and \allocOpen{} are open, in \allocSeconds{}\,s:

\begin{center}
\small
\begin{tabular}{lll}
\toprule
function & obligation & verdict \\
\midrule
\input{gen/alloc_table.tex}
\end{tabular}
\end{center}

The model test for the Python allocator notes that its model answers 0
past the end of an array where the compiled C reads out of bounds; the port
has no such read, and the obligations are the proof of that.

\subsection{Bounds, and the three that were false}\label{sec:bounds}

An earlier draft of this section reported four obligations open, all
\texttt{u64} additions, and said each fitted whenever the surrounding code
admitted it.  One did.  Three were false.

The first port of \texttt{bump} tested \texttt{used + n <= sizes[heap]}
--- evaluating \texttt{used + n} in order to find out whether it fitted.
At \texttt{used = n =} $2^{63}$ the sum wraps to 0, passes the test, and
\texttt{bump} returns 0; Crust's runtime check aborts on
\texttt{ensures used <= result}.  The kernel's proof of that
\texttt{ensures} was not wrong --- it is a theorem about \texttt{Nat}, where
nothing wraps --- and the open obligation was the place the model and the
machine part.  \texttt{slot\_addr} and \texttt{slot\_ok} formed
\texttt{bases[heap] + used} with nothing bounding either.  The first port
is kept in the paper's generator, and this run measures it again: the
addition in its guard is \oldGuardAdd{}, and the one behind the guard is
\oldReturnAdd{}.

Two things were missing to tell those apart.  The lift gave a slice's
elements no range, so nothing said \texttt{sizes[heap] <= u64::MAX}; it
now states \texttt{all\_le(sizes, max\_u64)} for every integer slice, as it
already stated a range for every integer parameter.  And
\texttt{by\_every\_bool} writes \texttt{true} or \texttt{false} in for a
guard and forgets that it held.  \texttt{by\_bounds} splits the same way,
but dependently --- the branch gets \texttt{Holds g} or
\texttt{g = false} as a hypothesis --- keeps the obligation's own
hypotheses rather than weakening them away, closes a branch whose facts
contradict each other, and at a leaf that does not compute chains
\texttt{<=} facts by transitivity.  Its steps are \boundsLemmas{} lemmas
added to the prelude, \boundsLemmasAll{} proved rather than assumed, among
them \texttt{add\_le\_of\_le\_sub}: \texttt{n <= s} and
\texttt{u <= s - n} give \texttt{u + n <= s}.  That is the lemma a
checked addition needs, and the fix to \texttt{bump} is its guard:

\begin{lstlisting}
if n <= sizes[heap] && used <= sizes[heap] - n { return used + n; }
\end{lstlisting}

The same answer wherever \texttt{used + n} fits; at $2^{63} + 2^{63}$ the
fixed build now answers \texttt{used}, $\overflowAnswer$.
\texttt{slot\_ok} answers 0 for an address past \texttt{u64::MAX}, since
such an address is in no region; \texttt{slot\_addr} cannot answer
anything but an address, so it asks its caller with
\texttt{\#[requires(bases[heap] <= u64::MAX - used)]}.  Writing that took
\texttt{u64::MAX} in Rust expressions, which Crust had flattened to an
undefined \texttt{u64\_MAX}; it is now a typed literal, and in the safety
lift the symbol \texttt{max\_u64} every range fact is stated against.

The same tactic closes \texttt{class\_covers} in \texttt{regs.rs}
(\classCoversProved{}), a gap the first draft pinned open: each arm of
\texttt{class\_for\_regs} is behind the range guard that bounds it, and the
\texttt{requires} bounds the last.

\section{Discussion: numerals that fit}
\label{sec:numerals}

Stating an overflow obligation needs the type's maximum, and the kernel
could not write one.  Its \texttt{Nat} is unary --- $3$ is
\texttt{succ (succ (succ zero))} --- so \texttt{u32::MAX} written out is
four billion nested terms, and the first attempt at stating one ran out of
memory.  \texttt{u64::MAX} is not writable at all.  The LeanOS paper
recorded the limit: ``a page-scale address is a term thousands deep''
\cite{hartshorn2026leanos}, and \texttt{crustproof} refused lengths past
200{,}000 for the same reason.

The kernel now has a literal node, as Lean's does.  It is the same term,
not an approximation: every closed numeral above zero normalises to one
literal, zero stays the constructor, so \texttt{succ (succ zero)} and the
literal $2$ share a normal form and are definitionally equal by the
ordinary comparison (this run: \litDefEq).  Where something needs a
numeral's structure --- the recursor matching on it, the unifier taking
\texttt{succ ?m} against it --- a literal is viewed as \texttt{succ} of the
numeral below it, one step at a time and no further.  The accelerated
definitions read and produce literals.  \texttt{numeral}$(2^{64}-1)$ is one
node (\litIsOneNode).

The effect on the SIMD contract certification of the earlier
paper~\cite{hartshorn2026memsafe}, which settles a bounds contract at a
known length:

\begin{center}
\begin{tabular}{lrl}
\toprule
length & seconds & verdict \\
\midrule
64 & \numSixtyFourSeconds & \numSixtyFourVerdict \\
$10^6$ & \numMillionSeconds & \numMillionVerdict \\
$2^{32}-1$ & \numUThirtyTwoSeconds & \numUThirtyTwoVerdict \\
$2^{64}-1$ (\texttt{len <= MAX}) & \numUSixtyFourSeconds & \numUSixtyFourVerdict \\
\bottomrule
\end{tabular}
\end{center}

and, so that speed is not mistaken for credulity, \texttt{div-by 4} at
$2^{32}-1$ is \numDivByFour.

Literal arithmetic is done by Python's integers, as Lean's is done by GMP
--- and two of the six 2026 Lean kernel bugs were in the interaction with
GMP~\cite{moses2026}.  The mitigation here is the one the earlier paper
introduced for the accelerators: each is run against the definition it
stands in for over a grid, in the prover's selftest, and a literal's
properties are pinned by \kernelChecks{} kernel checks (\kernelChecksVerdict)
and \hoareChecks{} prover checks (\hoareChecksVerdict).  It is a
mitigation, not a proof, and Section~\ref{sec:two} is the rest of the
answer.

One consequence of a literal being cheap arrived with \texttt{u64::MAX}
in a guard.  \texttt{sub 18446744073709551615 used} has nothing to
compute --- one argument is symbolic --- so the accelerator falls through
to the definition, and \texttt{sub} recurses on its first argument:
normalising it walks the literal down one \texttt{succ} at a time, under
the step function, $2^{64}$ times.  A literal above $2^{16}$ beside a
symbolic argument now leaves an accelerated definition folded.  That can
cost completeness and cannot cost soundness: reducing less can make the
kernel refuse a true statement, never admit a false one, and two copies of
the folded term are still equal.  Lean, given the same statements, checked
them without unfolding the literal either.

\section{Results: Independence, not superiority}
\label{sec:two}

In the summer of 2026, AI-assisted audits found soundness bugs in the Lean~4
kernel: Hulin's, Kumar's, and then six more in a systematic search by the
Lean FRO with OpenAI, after which the FRO changed its security practices
\cite{moses2026,reddit2026}.  Moses draws the conclusion in his title:
don't trust Lean alone.  His argument is not that Lean is bad --- it is that
a kernel is software, that its type theory is subtle, that its dependencies
are part of its trusted base, and that an agent given an impossible task
will look for the seams.

It would be easy, and wrong, to answer with a claim that our kernel is safer.
It is \kernelLines{} lines rather than eight thousand, and it can be read in
an afternoon; but small does not mean sound, and the people who missed Lean's
bugs were not careless.  The claim we make is about \emph{independence}.
Every theorem in this paper is checked twice: by a kernel written in Python
by one person for this purpose, and by Lean~4, written in C++ by others for
theirs.  The two share no code, no type theory implementation, no
arithmetic library.  A soundness bug lets a false theorem through only if it
exists in both, and a bug of the kind Hulin found --- an oversight in one
of three declaration paths --- has no reason to be repeated in a kernel
with a different shape.  The earlier ``A Second Kernel Agrees'' paper
\cite{hartshorn2026second} set this up for the LeanOS proofs; the Rust
theorems here go through the same export.  For the register sizer: Lean
\leanVersion{} \regsLeanVerdict{} the export in \regsLeanSeconds{}\,s, with
\regsLeanNoAxioms{} theorem depending on no axioms.

And not only the register sizer.  The prover keeps every obligation it
settles as a certificate --- the statement, the term \texttt{type\_check}
accepted, and the environment of callees, records and preconditions it was
proved in --- and \texttt{rustlean.py} (\rustleanLines{} lines) writes each
as a Lean file of its own, ending in \texttt{\#print axioms}.  One file per
theorem, because each obligation lives in its own environment and two may
define the same name differently.  Of the \rustCertificates{} certificates
from \texttt{regs.rs}, the safety table and \texttt{alloc.rs}, Lean accepted
\rustLeanAgreed{}, each depending on no axioms, in \rustLeanSeconds{}\,s.  A
theorem counts only if Lean accepts the file \emph{and} reports no axioms;
the open obligations have no certificate and are not counted either way.

Klaus et al.\ write that in their pipeline ``every proof is checked by the
Lean kernel, so AI output cannot compromise soundness''~\cite{klaus2026}.
That was the right thing to say in May 2026 and, after August, it is one
kernel's opinion.  Two kernels are cheap, and the second one being small
enough to read is what makes it worth having as a second.

\section{Conclusion: Why this dialect of Rust}

We call this, carefully, a dialect of Rust with logic safety on top of
memory safety: memory safety by Rust's rules, contracts checked at runtime
in every build, proofs from the shipped source, and independent re-checking,
in one compile.  To our knowledge no other Rust toolchain does all four at
once.

The reasons this is possible in Crust and not, or not easily, in
\texttt{rustc} are the reasons the series has given before.  The compiler
is \crustLines{} lines of Python that a reader can follow, so a new lift, a
new target, a new contract syntax is an afternoon rather than a project;
the same front end compiles C, C++, C\# and RPython into one translation
unit without an interface~\cite{hartshorn2026interop}, so a team can
sketch in Python, port to Rust, and regression-test both; the runtime
memory checks are placed where a proof could not remove
them~\cite{hartshorn2026memsafe}; and the operating system it targets was
written to be proved, with proofs that reach from the loader to the memory
map~\cite{hartshorn2026leanos}.  This paper adds the layer those left open.

What this is not:

\begin{itemize}
\item \textbf{Not a proof about machine code.}  The theorems are about the
  lifted model; the model is tied to the binary by tests that run both over
  the same inputs, and a corpus can miss an input.  Translation validation
  from IL to emitted code is the far end of this road, as before.
\item \textbf{Arithmetic, but only chains of it.}  \texttt{by\_bounds}
  chains \texttt{<=} facts through addition, a checked subtraction and a
  slice's range; \texttt{a < 1000} and \texttt{b < 1000} give
  \texttt{a + b < 1999} only by \texttt{add\_le\_add} and a literal
  comparison, and nothing about $\times$ is proved yet.  Its search is
  bounded in depth, so an open obligation still means only that it was not
  found.
\item \textbf{Not signed.}  Signed integers are refused, not modelled.
\item \textbf{Not all of Rust.}  Methods, traits, closures in the proof
  subset, recursion: refused by name.  The compiler handles far more Rust
  than the lift does; the lift grows toward it.
\item \textbf{Known gaps, asserted.}  A record contract through a field
  update is not yet proved; a test asserts it is still open, so closing it
  cannot happen quietly, and it is checked at runtime meanwhile.  A
  callee's \texttt{requires} that mentions \texttt{u64::MAX} reaches its
  callers as the number, not the symbol, so a caller cannot yet discharge
  it.
\end{itemize}

\appendix
\footnotesize
\section{Reproducing this paper}

\texttt{python3 rustproof\_paper.py} in the RosettaMath tree, beside a
checkout of Crust, writes every number above into
\texttt{gen/rustproof\_facts.tex} and the lifted models and Lean sources
into \texttt{gen/}; \texttt{pdflatex rustproof.tex} typesets it.  The
tables are generated; the listings are the shipped files.  In the Crust
tree, \texttt{make test\_rustproof} runs the \rustproofTests{} tests behind
Sections~\ref{sec:lift}--\ref{sec:safety}, and \texttt{make prove\_rust}
prints the obligation report for every Rust module under \texttt{leanos/}.


\begin{thebibliography}{9}
\bibitem{moses2026} Milo Moses.  ``Don't trust Lean4 alone.''  LessWrong,
  17 September 2026.
  \url{https://www.lesswrong.com/posts/jgmmMa7AqJNausrqx/don-t-trust-lean4-alone}
\bibitem{reddit2026} ``Model by OpenAI discovers 6 more soundness bugs in
  the Lean kernel.''  r/mathematics, August 2026.
  \url{https://www.reddit.com/r/mathematics/comments/1vdb5an/}
\bibitem{biggar2020} Paul Biggar.  ``First thoughts on Rust vs OCaml.''
  Darklang blog, August 2020.
  \url{https://blog.darklang.com/first-thoughts-on-rust-vs-ocaml/}
\bibitem{klaus2026} Natalia Klaus, Juan Conejero, Palina Tolmach.  ``A
  Rust-to-Lean Verification Pipeline with AI Provers: An Experience
  Report.''  arXiv:2605.30106, 2026 (AIMACS at CAV 2026).
  \url{https://arxiv.org/abs/2605.30106}
\bibitem{creusot} Xavier Denis, Jacques-Henri Jourdan, Claude March\'e.
  ``Creusot: a Foundry for the Deductive Verification of Rust Programs.''
  ICFEM 2022.
\bibitem{prusti} Vytautas Astrauskas, Peter M\"uller, Federico Poli,
  Alexander J. Summers.  ``Leveraging Rust Types for Modular Specification
  and Verification.''  OOPSLA 2019.
\bibitem{hartshorn2026leanos} Brent S. Hartshorn.  ``LeanOS: a kernel
  written to be proved by LEAN4, from a proof-checked loader to a
  proof-checked memory map.''  SSRN, 2026.
  \url{https://dx.doi.org/10.2139/ssrn.7461558}
\bibitem{hartshorn2026lean} Brent S. Hartshorn.  ``LEAN Proofs Small Enough
  to Read: Kernel Contracts from \LaTeX{} Theorems to Compiler Decisions.''
  ai.viXra.org:2609.0019, 2026.  \url{https://ai.vixra.org/abs/2609.0019}
\bibitem{hartshorn2026second} Brent S. Hartshorn.  ``A Second Kernel
  Agrees: `LEAN Proofs Small Enough to Read', Checked by Lean4 and Read Four
  Ways.''  SSRN, 2026.  \url{https://dx.doi.org/10.2139/ssrn.7443439}
\bibitem{hartshorn2026memsafe} Brent S. Hartshorn.  ``Memory Safety Where
  it is Needed: Proof-guided Runtime Checking in a Toolchain Small Enough to
  Read.''  SSRN, 2026.  \url{https://dx.doi.org/10.2139/ssrn.7396160}
\bibitem{hartshorn2026interop} Brent S. Hartshorn.  ``Interoperation
  Without an Interface: C++ and Rust in One Translation Unit, and One
  Toolchain.''  SSRN, 2026.  \url{https://dx.doi.org/10.2139/ssrn.7315678}
\end{thebibliography}

\end{document}
